%========================================================
% CHAPTER 4 — Experimental Activity and Polarization
%             Control in Fiber-Based Quantum Links
%========================================================

This chapter presents the experimental activity carried out at the POLICOM
laboratory of the Department of Electronics, Information and Bioengineering
(DEIB) at Politecnico di Milano.


%========================================================
% SECTION 4.1 — Experimental Objectives and Development
%               Strategy
%========================================================

\section{Experimental Objectives and Development Strategy}
\label{sec:experimental_objectives}

The main objective of the experimental work is to investigate, characterize, and
control the polarization behavior of light transmitted through a two-arm
fiber-optic system, with particular attention to the mechanisms that may affect
the transmission of polarization-encoded quantum states. The experimental
activity therefore complements the analytical and numerical framework developed
in the previous chapters by providing a laboratory-based study of the same
physical phenomena.

The adopted architecture is intentionally simplified. Rather than implementing
a complete quantum-communication protocol, the setup reproduces a representative
subsystem of a future entanglement-based BBM92 architecture. In this
context, the objective is to study the behavior of the optical components and
polarization-control stages that would be required in a realistic fiber-based
installation, and to identify the experimentally accessible observables that can
be used to monitor their operation.

The experimental platform combines several optical and electronic components
previously characterized and calibrated in the laboratory, including a
spontaneous-parametric-down-conversion source, polarization-control elements,
polarization beam splitters, polarization-maintaining fiber sections,
single-photon avalanche detectors, time-tagging electronics, and, in the later
development stages, an electronically controlled polarization controller. These
components are progressively integrated into a common measurement and control
framework.

The experimental activity was developed according to three main stages:

\begin{enumerate}
    \item \textbf{Manual polarization control.}
    The first stage establishes the basic two-arm detection architecture using
    four single-photon avalanche detectors and manually adjusted polarization
    control. This configuration provides the reference measurement platform for
    the characterization of single-photon count rates and coincidence
    observables.

    \item \textbf{Automatic polarization control.}
    The second stage introduces electronic polarization control and
    the corresponding software-based feedback system. The controller is first
    calibrated and then used to investigate automatic polarization alignment.
    Different feedback and optimization strategies are tested in order to
    evaluate their ability to recover and maintain a desired polarization
    condition from experimentally measured coincidence observables.

    \item \textbf{Controlled PM-fiber propagation experiment.}
    The final stage combines the four-detector analysis, automatic polarization
    control, and polarization-maintaining (PM) fiber sections in a single
    experimental setup. The PM fiber is used to introduce a controlled
    birefringent propagation mechanism and to investigate the resulting
    polarization degradation under experimentally reproducible conditions.
\end{enumerate}

The final PM-fiber experiment is performed using a non-entangled input polarization state. 
Consequently, it does not directly measure entanglement degradation, 
but it allows the polarization-decoherence and depolarization mechanisms 
induced by controlled birefringent propagation to be experimentally characterized.
The measured coincidence observables are then interpreted
through the calibrated physical PM-fiber model developed in the previous
chapters.

More generally, the scope of the experimental activity is not to reconstruct the
microscopic evolution of each individual photon inside the fiber. The objective
is instead to determine how the available optical components, polarization
controllers, detectors, and fiber sections behave when assembled into a realistic
quantum-communication link. In this sense, the laboratory platform provides an
engineering-level validation environment connecting the theoretical fiber models
with experimentally accessible quantities and with the requirements of a future
BBM92-oriented implementation.

%========================================================
% SECTION 4.2 — Experimental Platform and Instrumentation
%========================================================

\section{Experimental Platform and Instrumentation}
\label{sec:experimental_platform}

This section describes the common hardware and acquisition infrastructure used
throughout the experimental activity.

The presentation includes the SPDC source and the optical architecture inherited
from the previously developed laboratory platform, manual polarization
controllers, the electronic polarization controller, polarizing beam splitters,
polarization-maintaining fibers, SPAD detectors, and the time-tagging system.

The software developed for detector configuration, real-time singles monitoring,
coincidence-histogram acquisition, timing-delay compensation, EPC control, and
experimental data storage is also introduced.

The description focuses on the role played by each component in the experiments
performed in this work, while detailed characterization of instrumentation already
documented in previous laboratory work is only recalled where necessary.


%========================================================
% SECTION 4.3 — Phase I: Manual Polarization Control and
%               Four-SPAD Measurements
%========================================================

\section{Phase I: Manual Polarization Control and Four-SPAD Measurements}
\label{sec:phase1_manual_control}

The first experimental phase establishes the reference measurement configuration
using four SPAD detectors and two manually operated polarization controllers.

The purpose of this stage is to validate the complete acquisition chain, including
single-photon count-rate measurements, coincidence-peak identification, timing
alignment between detector channels, and polarization-dependent redistribution
of coincidence events.

Manual polarization control is used to align the polarization states with the
measurement basis and to investigate how changes in the optical configuration
modify the observed coincidence channels.

This phase provides the experimental baseline required before introducing
automatic polarization compensation.


%========================================================
% SECTION 4.4 — Phase II: Automatic Electronic
%               Polarization Control
%========================================================

\section{Phase II: Automatic Electronic Polarization Control}
\label{sec:phase2_epc}

The second phase introduces an electronic polarization controller into a reduced
experimental configuration without PM fiber.

One photon provides the reference detection event, while the polarization of the
second optical arm is controlled by the EPC before being analyzed by a PBS and
two SPAD detectors.

The polarization-compensation problem is formulated as a closed-loop,
derivative-free optimization problem in which the EPC control voltages are
updated according to experimentally measured coincidence observables.

This section introduces the feedback architecture, the EPC control variables,
the measurement cycle, and the general optimization framework used throughout
the automatic-control experiments.


%========================================================
% SECTION 4.5 — Coincidence-Based Objective Functions
%========================================================

\section{Coincidence-Based Objective Functions}
\label{sec:epc_objective_functions}

This section defines the experimental quantities used to provide feedback to the
electronic polarization controller.

The principal observables are the net coincidence counts associated with the two
outputs of the polarization analyzer. These quantities are extracted from the
measured correlation histograms through coincidence-peak localization,
integration, and background estimation.

The main objective functions used in the experiments are introduced and compared,
including optimization of a selected coincidence channel and combined objective
functions constructed from both polarization outputs.

Particular attention is given to the physical meaning of the objective functions
and to the distinction between single-photon count rates and coincidence-based
feedback.


%========================================================
% SECTION 4.6 — EPC Optimization Algorithms
%========================================================

\section{EPC Optimization Algorithms}
\label{sec:epc_optimization_algorithms}

This section presents the principal optimization algorithms implemented for
automatic EPC control.

The first method is sequential coordinate search, in which the EPC control
parameters are optimized one at a time while retaining configurations that
improve the measured objective function.

The second method is a stochastic coarse-to-fine optimization based on simulated
annealing. A coarse exploration of the EPC voltage space is followed by a finer
local search, while probabilistic acceptance of non-improving configurations
provides a mechanism for escaping unfavorable local solutions.

The algorithms are discussed in terms of convergence behavior, measurement cost,
repeatability, sensitivity to photon-counting fluctuations, and suitability for
closed-loop polarization control.


%========================================================
% SECTION 4.7 — Phase III: Controlled Propagation Through
%               PM Fiber
%========================================================

\section{Phase III: Controlled Propagation Through PM Fiber}
\label{sec:phase3_pm_fiber}

The final experimental phase introduces a controlled polarization-maintaining
fiber section into the optical path.

The experimental configuration is

\begin{equation}
\text{SPDC OUT 1}
\rightarrow
\text{Manual PC}
\rightarrow
\text{PBS}_{45^\circ}
\rightarrow
\text{PM}_{45^\circ}
\rightarrow
[\text{PM}_{0^\circ}]^{N}
\rightarrow
\text{EPC}
\rightarrow
\text{PBS}
\rightarrow
\{\text{SPAD 3},\text{SPAD 4}\},
\end{equation}

while the reference arm is

\begin{equation}
\text{SPDC OUT 2}
\rightarrow
\text{SPAD 1}.
\end{equation}

Unlike the entanglement-distribution models developed in the theoretical part of
the thesis, this experiment uses a non-entangled input polarization state.

Its purpose is to experimentally isolate and characterize the effect of
controlled birefringent propagation and differential group delay on the measured
polarization observables.


%========================================================
% SECTION 4.8 — Role of the Rotated and Aligned PM-Fiber
%               Sections
%========================================================

\section{Role of the Rotated and Aligned PM-Fiber Sections}
\label{sec:pm_fiber_configuration}

This section explains the physical motivation for the sequence composed of an
initial PM-fiber section rotated by approximately \(45^\circ\), followed by
additional aligned PM-fiber sections.

The rotated section couples the prepared input polarization to both principal
polarization modes of the PM fiber, allowing the fast and slow components to
propagate simultaneously.

The subsequent aligned PM-fiber sections increase the accumulated differential
group delay while preserving, to first approximation, the same principal-axis
orientation.

The number of aligned sections, or equivalently the total PM-fiber length, is
therefore used as the main experimental control parameter governing the strength
of the polarization--temporal separation.


%========================================================
% SECTION 4.9 — Measured Coincidence Observables Under
%               PM-Fiber Propagation
%========================================================

\section{Measured Coincidence Observables Under PM-Fiber Propagation}
\label{sec:pm_fiber_observables}

This section presents the quantities directly accessible with the current
experimental instrumentation.

Since the available setup performs polarization analysis in a single measurement
basis, full quantum-state tomography is not carried out. Consequently, quantities
such as the complete density operator, concurrence, and state fidelity are not
directly reconstructed from the experimental measurements.

The experimental analysis instead focuses on the coincidence observables
associated with the two PBS output channels and on their evolution as the
PM-fiber configuration is modified.

The measurements are used to quantify the redistribution and loss of polarization
contrast produced by increasing differential propagation between the principal
fiber modes.


%========================================================
% SECTION 4.10 — Comparison with a Calibrated Physical
%                PM-Fiber Model
%========================================================

\section{Comparison with a Calibrated Physical PM-Fiber Model}
\label{sec:pm_fiber_model_comparison}

The experimental results are interpreted through a calibrated physical model of
the PM-fiber propagation rather than through complete quantum-state tomography.

The model incorporates the relevant properties of the PM-fiber sections, such as
their length, principal-axis orientation, birefringence, differential group delay,
and the spectral characteristics of the propagating photons.

From these parameters, the expected polarization evolution and the corresponding
measurement probabilities at the PBS outputs are calculated and compared with
the experimentally measured coincidence observables.

This model-assisted approach allows the physical origin of the observed
polarization degradation to be investigated within the limitations imposed by
the available measurement basis.


%========================================================
% SECTION 4.11 — Connection with the Theoretical and
%                Numerical Framework
%========================================================

\section{Connection with the Theoretical and Numerical Framework}
\label{sec:experimental_theoretical_connection}

The final section connects the experimental observations with the physical
fiber-propagation framework developed in the previous chapters.

The PM-fiber experiment provides a controlled realization of
frequency-dependent birefringent propagation, in which polarization becomes
coupled to temporal or spectral degrees of freedom through differential group
delay.

When these additional degrees of freedom are not resolved by the measurement,
the resulting polarization dynamics can be represented through an effective
reduced-state description.

The experiment therefore provides a physical test of the propagation mechanism
underlying the numerical PMD model. The distinction between theory and experiment
is maintained explicitly: the numerical framework treats both single-photon and
entangled two-photon states, whereas the final laboratory experiment investigates
the same physical mechanism using a non-entangled input state and a restricted
set of coincidence observables.